\section{Discussion}
\label{sec:discussion}

\subsection{Interpretation of the findings}

The strongest result is not that \routeB\ produces a more attractive square
texture. Its flat atlas is often less legible than \routeA\ because it stores
pieces of a continuous composition. The relevant result is that the composition
reconstructs on the weapon while satisfying a true asset-seam constraint. This
changes the unit of evaluation from the generated rectangle to the mapped asset.

The paired experiment supports that representation choice for the tested AK-47
candidates. All four \routeB\ candidates improve both seam and multi-view
performance. The seam gain is expected in part because \routeB\ explicitly
contains an asset-edge operation, but the simultaneous multi-view gain is
important: continuity was not purchased by simply flattening all visible detail.
The separate before/after table further isolates the seam correction and shows
that the operation, rather than the generator alone, is responsible for the low
final error.

The live case complicates a simple ``higher score is better'' interpretation.
\routeA\ obtains a higher multi-view score than the selected \routeB\ texture, yet
fails the asset-seam threshold by a large margin. This is precisely the situation
in which feasibility must precede preference. The edge-width ablation reaches the
same conclusion under a fixed candidate pool: a small gain in multi-view score
causes the weighted policy to choose an invalid candidate.

The rollback results are also informative. None of the formal \routeC\ attempts
improves the locked score, and the live correction is distinctly worse. Reporting
these outcomes avoids presenting the Agent as an automatically self-improving
system. The contribution is narrower and more defensible: the Agent can identify a
target, produce local alternatives, prove locality, and decline to replace a
better baseline.

\subsection{Metric appraisal}

The evaluation stack deliberately contains several measures because no single
metric captures deployable visual quality. Asset-seam error has a clear technical
interpretation, but a very low value can be achieved by over-smoothing boundaries.
Multi-view retention protects against that failure, although it still summarizes
contrast, saturation, detail, coverage, and consistency rather than aesthetic
quality.

The component-detail balance metric failed on high-contrast weapon-space
candidates: several \routeB\ values collapsed to zero and made the locked Agent
score lower even when seam and multi-view performance improved. For this reason,
the component-balance and aggregate-score differences are not used as evidence
that \routeA\ is superior. They are retained as a documented metric limitation and
as motivation for reporting primary measures separately.

Mapped-element readability has a similar boundary. The selected artwork scores
highly at the source and substantially lower on the top view, which is useful
diagnostic information. However, it depends on model-based semantic judgments and
view weighting chosen for this case study. The human artwork decision remains
authoritative.

\subsection{Internal and external validity}

Several controls strengthen internal validity: fixed assets and cameras, shared
seed schedules within the formal comparison, exact source reuse after human
selection, machine-readable logs, isolated sub-ablations, and deterministic
rendering. The renderer, UV bake, seam metric, selection policy, and locality gate
also have regression tests.

The main threat is sample size. Four formal pairs are sufficient to expose
consistent behavior in the preserved candidate set but not to estimate performance
over the space of themes, generators, or weapons. Route-specific generation logic
also means the comparison evaluates two complete design strategies rather than a
single-factor image transformation. This is appropriate for the research question,
but the result should not be interpreted as a causal estimate of one isolated
algorithmic component.

External validity is limited by one complete weapon, one principal creative case,
one software renderer, and a small set of backends. The M4A4 result is useful
because it fails the joint threshold: it shows that adapter portability is not the
same as parameter portability. Broader claims would require several weapon
families, game material systems, viewing conditions, and users.

\subsection{Human-Agent interaction limitations}

The three checkpoints solve a structural problem: they separate interpretation,
creative direction, and mapped artwork. They also make disagreement recoverable.
A user can edit a theme before generation, reject an unwanted direction, or reject
an artwork after mapping without discarding the complete run.

The present study verifies this contract technically, but does not compare it with
a one-prompt baseline through a user study. Future work should measure task time,
number of revisions, perceived control, preference confidence, and the frequency
with which source-only and mapped choices disagree. Such a study would provide
stronger evidence for RQ4 beyond successful execution.

\subsection{Material-channel and engine-validation boundary}

The implemented output is a base-colour texture for the Custom Paint Job
workflow. It does not generate aligned normal, roughness, metallic, height, or
displacement maps, and it should therefore not be described as a complete
physically based material authoring system. The current contribution concerns
semantic planning, colour-texture composition, UV-aware baking, mapped-view
evaluation, and constrained selection.

CS2 Workbench captures provide important deployment evidence, but their lighting
and reflectance are produced by the engine and finish configuration. They do not
demonstrate model-generated surface relief or material response. The known-good
calibration case verifies the UV orientation and import contract, while the
black-and-gold marble case verifies coherent in-engine colour-texture deployment.
These cases are kept distinct from the accepted live dragon Agent run, whose
final TGA still requires a fixed left/right/top Workbench capture for the strongest
end-to-end visual record.

The most direct technical extension is a multi-channel material pipeline that
generates spatially aligned base-colour, normal, roughness, and metallic maps,
checks cross-channel consistency, and validates the result under fixed engine
lighting. This would require material-aware metrics and export adapters in
addition to the present UV and colour-texture checks.

\subsection{Project management and methodological evolution}

The registered topic, \emph{Latent Vision $\rightarrow$ Image Generation}, was
intentionally broad and emphasized image-to-image workflows, diffusion models,
conditioning, and creative visual outputs. Supervisor guidance later superseded
that original scope: given the remaining delivery window, completing the technical
development of an image-processing Agent was considered sufficient, and the Agent
did not need to preserve the original Latent Vision framing. The team therefore
treated SkinSmith as a focused implementation of the approved image-processing
Agent direction rather than as an attempt to cover every objective in the original
topic description.

Work then progressed through dated milestones: deterministic generation and seam
repair; real OBJ rendering; local diffusion; UV and seam diagnostics; weapon-space
composition; formal A/B/C experiments; provider-backed generation; the reusable
Agent runtime; the three-checkpoint client; and academic production.

Several changes were driven by failed evidence rather than by the original plan.
Direct border blending was removed after it produced a visible mirrored frame.
An early mask-aware route was retained as negative evidence because it changed
component appearance without creating a coherent whole-weapon composition.
Per-component dragon generation was abandoned after it produced semantically wrong
source parts and framed panels. A legacy UV assumption was rejected after
deployment mismatch. These corrections delayed the interface, but reduced the risk
of polishing a workflow whose asset contract was wrong.

The group used five connected work packages with one integration lead. YUAN Ye
served as project lead and integration owner, defining the system architecture and
coordinating the shared Agent contract. CHEN Yuhong owned creative planning and
literature synthesis; Ben Tangjie owned generation-backend integration and
reproducibility; Wang Zhengye owned asset, UV, and rendering validation; and WANG
Lihui owned evaluation, experiment reporting, and submission quality control.
Each work package contributed method decisions, implementation or configuration,
experimental verification, and documentation to the final system. This structure
gave every member an assessable technical contribution while retaining one owner
for cross-module consistency.

\subsection{Responsible use, provenance, and rights}

The project uses generative models as experimental components. Structured theme
and style planning in the accepted live case used
\texttt{gemini-3.1-flash-lite}; source artwork used
\texttt{gemini-3.1-flash-image}; and SD-Turbo provided a pinned local baseline.
Prompts, model identifiers, validation traces, and source hashes are recorded.
Credentials are supplied through local environment variables and removed from
evidence.

The system is restricted to original themes and explicitly prohibits copied skins,
logos, brands, teams, protected characters, text, watermarks, and imitation of a
named living artist. Generated sources are reviewed before mapping. Valve geometry
is used only as an external case-study asset and is not redistributed in the
repository. The report does not claim ownership of Valve assets, affiliation with
Valve, or readiness for Steam Workshop publication.

These safeguards do not resolve every concern. The training data of external image
models are not controlled by this project, and automated filters can miss
similarity or embedded content. Any use beyond the academic prototype would
require model-licence review, independent rights checking, stronger moderation,
game-specific terms review, and a clear appeals process for rejected or disputed
content.

\subsection{Practical implications}

Three practical lessons follow from the study:

\begin{enumerate}
  \item A source image should not be approved before it is mapped. Candidate cards
  should bind the source to the views on which it will be judged.
  \item A deployment constraint should not be averaged into a general quality
  score. Feasibility and preference answer different questions.
  \item Agent refinement should have a budget and a rollback rule. In creative
  pipelines, a plausible correction can still remove valuable detail.
\end{enumerate}

These lessons extend beyond weapon skins to other assets whose appearance is
stored in fragmented atlases, including props, clothing, vehicle parts, and
product visualizations. The geometry and deployment adapters would change, but
the evidence and decision structure would remain relevant.
